% ===========================================================================
% SOURCES: docs/claude_project/04_VERIFICATION.md
%          docs/references/rtl_review.md (the 43 findings)
%          docs/claude_project/07_RTL_REVIEW_STATUS.md
%          sim/** (31 testbenches)
%          docs/lab_notebook.md, 2026-10-07 (the first board run - what the
%            method bought when it met silicon)
%
% THIS CHAPTER IS THE THESIS'S METHODOLOGICAL CONTRIBUTION. It was §6.2 in the
% first outline - four subsections inside the design chapter - which badly
% understated it. Budget 12-15 pages.
%
% The claim to defend: bit-exactness is not a testing strategy, it is a design
% method. It changes what questions can be asked. Write it so a committee
% member who builds accelerators recognises something they do not do.
% ===========================================================================
\chapter{Verification Methodology}
\label{chap:verification}

\section{The Contract: Bit-Exactness Against Golden Vectors}
\label{sec:bit-exact-contract}
% The software model emits the exact integer output of every layer; every RTL
% block and the assembled network reproduces those numbers bit for bit.
% Four consequences, each worth a paragraph:
%   1. there is never an argument about "close enough" - it matches or it is
%      a bug;
%   2. a failure localises to one layer instead of to "the network";
%   3. the hardware-vs-software comparison in Chapter~\ref{chap:evaluation} is not a
%      separate experiment, it IS the test suite;
%   4. every transformation in Chapter~\ref{chap:timing} was provably a reassociation
%      or a latency change and never a change of value, so one \thElements-element
%      check regresses the entire timing campaign.
% Point 4 is the one that made a 43-finding review executable in two days.
% Say so explicitly - it is the strongest practical argument in the chapter.

\section{Three Levels of Verification}
\label{sec:three-levels}
% unit -> integration -> system, and the important claim: each level found
% bugs the level below STRUCTURALLY could not.
%   - unit: one engine against a golden vector, no controller;
%   - integration: a real layer, windows and addresses built by hand in the
%     testbench - which is why those testbenches became the controller's
%     specification;
%   - system: accel_tb, the whole network, \thElements elements over
%     \thInstructions instructions.
% \DD{017} is the evidence: two busy bugs that no standalone testbench could
% have found, which appeared the moment out_stage was shared.
% Note also that each level was added BEFORE the phase that needed it.

\section{Mutation Testing}
\label{sec:mutation}
% How do you know a passing testbench tests anything? Inject a fault and
% require the suite to fail. pe_array was mutated three ways on 2026-09-11,
% and the directed cases were designed so that random data could not hide a
% symmetric wiring mistake.
% This is the section that answers "your tests pass, but are they any good?" -
% a question a committee will ask and most theses cannot answer.

\section{Simulation Assertions}
\label{sec:assertions}
% The assertion layer catches what golden vectors structurally cannot: a
% handshake violated, a FIFO overrun, an address out of range - failures that
% do not change the final numbers on the passing path. Give an example where
% an assertion fired on something the element comparison would have missed.

\section{Guards Against Stale Artefacts}
\label{sec:guards}
% The \DD{019} aftermath, as process rather than as arithmetic
% (Section~\ref{sec:acc-bound} has the arithmetic). Four guards, each closing a
% failure that actually occurred:
%   - export.py recomputes the accumulator bound and exits non-zero if the
%     network's worst case does not fit the ACC_W it reads out of the RTL;
%   - gen_program.py refuses a checkpoint-hash mismatch and a MISSING hash -
%     the missing case is the one that bit, because the export in git was five
%     weeks older than the checkpoint on disk and neither artefact said so;
%   - gen_program.py asserts the residual M <= 1 property;
%   - the layer testbenches' per-tensor scalars are GENERATED
%     (hardware/tb/model_scalars.svh, by gen_program.py --emit). Found by the
%     redeploy of \DD{036}: four testbenches had the old model's avgpool and
%     features.3.add multipliers typed in, and three of them failed - 288 of
%     1,280 avgpool channels a code or two low - while the RTL was unchanged
%     and accel_tb was bit-exact. (The fourth had stopped on a file lock
%     first; it was found by reading it.)
%     The suite caught its own stale data, which is the method working, and
%     the fix made it impossible rather than corrected it.
% The general statement: the project's defence against stale numbers is that
% the build fails, not that someone remembers. The thesis's own numbers follow
% the same rule (Section~\ref{sec:repro-numbers}): the redeploy moved every
% accumulator figure, and the ones still copied by hand were the ones that
% went wrong.

\section{The External RTL Review}
\label{sec:rtl-review}
% 43 findings, and what happened to each (the full table belongs in
% Chapter~\ref{appendix:decisions} or an appendix, not inline). What to draw out here:
% several findings were closed DIFFERENTLY from how the review proposed,
% because "is this number measured or bounded?" gave a better answer than the
% recommendation - \DD{022} is the clearest case, where the path was decomposed
% before the code and one register replaced a recommended grouped chain.
% Be honest about the still-open ones (C-10, C-11, T-6, T-8, T-9, \DD{023});
% Section~\ref{sec:future} carries them.

\section{What This Method Does Not Cover}
\label{sec:verification-limits}
% Required for honesty, and it pre-empts the obvious attack:
%   - until \DD{033} the \thElements-element check started at 000_input.hex,
%     so resize and quantisation were unverified (Section~\ref{sec:preproc} closed it);
%   - bit-exactness is checked on \thImagesVerified{} images, one per class, not on
%     the test set - it is a check of the ARITHMETIC, and
%     Section~\ref{sec:accuracy} is what measures the model. Since DD-036's
%     redeploy those images are test-split photographs the model never saw -
%     the first of each class in file order that it classifies correctly;
%     before it they were training or validation images;
%   - simulation is not silicon - and that line now has an answer: on
%     2026-10-07 the board reproduced a golden image's four logits exactly
%     (Section~\ref{sec:board-remaining}), and on 2026-10-22, under Linux, all
%     four golden images, one per class, through the RGB path
%     (Section~\ref{sec:camera-path}). The method is WHY each was a comparison
%     rather than a debugging session: the board had an exact answer to
%     match, not a plausible-looking class. Say what it does not cover: four
%     images is a check of the arithmetic on silicon, not of the model.
